\section{周期结构如何补上缺少的动量}
\label{chap:feqo-structured-media}

令结构沿电子轨迹方向的周期为 \(a\)，其倒格矢为 \(G_m=2\pi m/a\)。场的一个 Bloch 分量写作 \(E_m(z,\omega)=u_m e^{\ii(k_z+G_m)z}\)。它与运动电子在\cref{eq:feqo-beta} 中的相位相乘，并在长度 \(L\) 的对称作用区积分，得到
\begin{equation}
\int_{-L/2}^{L/2}\dd z\,
e^{\ii(k_z+G_m-\omega/v)z}
=L\sinc\!\left[\frac{(k_z+G_m-\omega/v)L}{2}\right].
\label{eq:feqo-periodic-phase-matching}
\end{equation}
长结构使相位匹配窗口收窄到 \(|k_z+G_m-\omega/v|\lesssim1/L\)。光栅给或取整数个 \(\hbar G_m\) 后，原先落在光锥外的电子近场可能耦合到传播光子；但 \(u_m=0\) 时积分仍为零，说明相位匹配只给允许条件，强度还取决于模式场与电流的重叠。

为什么非要周期结构不可，看一张几何图就清楚。\cref{fig:light-cone-synchronism} 把光锥 \(\omega/c=|k_z|\) 与电子线 \(\omega/c=\beta_0|k_z|\) 画在同一个 \((k_z,\omega)\) 平面上。取 \(\beta_0=0.70\)，两条直线只在原点相遇，其余各处中间的阴影区就是 \(\omega/v=k_z>\omega/c\) 的部分：电子一路能提供的纵向波数全都落在光锥以外，因此在平移对称的理想几何里它无法辐射。要让 \(k_z+G_m\) 落进光锥，只能在横轴上平移整数个倒格矢，这就是\cref{eq:feqo-periodic-phase-matching} 里必须有非零 \(u_m\) 的几何含义。

\begin{figure}[htbp]
\centering
\includegraphics[width=0.72\linewidth]{light_cone_synchronism.pdf}
\caption{光锥与电子线的相对位置，取 \(\beta_0=0.70\)。横轴为 \(k_z/k_0\)，纵轴为 \(\omega/(ck_0)\)；实线是光锥 \(\omega/c=|k_z|\)，虚线是电子线 \(\omega/c=\beta_0|k_z|\)，阴影区是电子能提供的波数全部落在光锥之外的区域，只能靠倒格矢平移进入。}
\label{fig:light-cone-synchronism}
\end{figure}

例如出射光在真空中与电子轨迹夹角为 \(\theta\)，其纵向波数为 \(k_z=(\omega/c)\cos\theta\)。把它代入相位匹配式，取 \(m>0\) 与 \(\lambda=2\pi c/\omega\)，得
\begin{equation}
\lambda=\frac{a}{m}
\left(\frac{c}{v}-\cos\theta\right).
\label{eq:feqo-smith-purcell-dispersion}
\end{equation}
这条关系说明电子速度、光栅周期和收集角共同决定可见波长。由于 \(v<c\)，括号在任何实角都为正；改变 \(a\) 会按比例移动谱线，而改变探测角会在同一结构上选不同波长。有限作用长度时并不存在严格的单一角度：\cref{eq:feqo-periodic-phase-matching} 的第一零点满足 \(|\Delta k|=2\pi/L\)，在固定 \(\omega\) 下对应 \(|\Delta\cos\theta|\simeq\lambda/L\)。这个估计只给相位匹配宽度，实际角分布还乘模式外耦合、颗粒尺寸及收集效率。若无相应的 \(u_m\)，即使满足\cref{eq:feqo-smith-purcell-dispersion} 也没有这一级辐射。

材料中的两个裸模也可相互混合。先在无损、只含两个归一化模的模型中定义
\[
H_{\rm mix}=\hbar\begin{pmatrix}\omega_a&g\\g^*&\omega_b\end{pmatrix}.
\]
解二阶特征方程得 \(\omega_\pm=(\omega_a+\omega_b)/2\pm\sqrt{(\omega_a-\omega_b)^2/4+|g|^2}\)。当 \(\omega_a=\omega_b\) 时分裂为 \(2|g|\)；本征向量决定电子电流对上下两支各有多大重叠。给裸模另加衰减率 \(\kappa_a,\kappa_b\) 后，有效矩阵的对角元可写成 \(\omega_{a,b}-\ii\kappa_{a,b}/2\)，但它不是封闭系统的 厄米 哈密顿量；复本征值描述峰位和线宽的近似，而峰是否能被看成两条，还取决于耦合、线宽、探测窗口和背景干涉。

上面的 \(\omega_\pm\) 在无量纲坐标里就是一对双曲线。\cref{fig:hopfield-avoided-crossing} 的横轴是失谐 \((\bar\omega_c-\omega_m)/|g_{\rm cm}|\)，纵轴是 \((\omega-\bar\omega)/|g_{\rm cm}|\)，两条灰线是未混合的重整化类光子模与裸物质模，在原点交叉；两条实线是混合后的两支，在失谐为零处不相交，而是一上一下停在 \(\pm1\)，即最小劈裂正好等于 \(2|g_{\rm cm}|\)。失谐推到 \(\pm8\) 时实线几乎贴在虚线上，说明远离共振时两个裸模各自复位；可分辨的劈裂只存在于 \(|\bar\omega_c-\omega_m|\) 与 \(|g_{\rm cm}|\) 同量级的窗口内。

\begin{figure}[htbp]
\centering
\includegraphics[width=0.76\linewidth]{hopfield_avoided_crossing.pdf}
\caption{RWA 下的反交叉。横轴为重整化失谐 \((\bar\omega_c-\omega_m)/|g_{\rm cm}|\)，纵轴为相对频率 \((\omega-\bar\omega)/|g_{\rm cm}|\)；虚线为重整化类光子模，点线为裸物质模，两条实线为混合后的本征频率 \(\omega_\pm\)，在 \(\bar\omega_c=\omega_m\) 处的最小劈裂为 \(2|g_{\rm cm}|\)。}
\label{fig:hopfield-avoided-crossing}
\end{figure}

把运动学与收集通道放在同一实验里，电子可激发满足\cref{eq:feqo-periodic-phase-matching} 的内部模；材料损耗决定其中多少变热；边缘或光栅耦合决定多少成为辐射；最后光学系统只读收集到的角度与偏振。这几个比例分别变化，所以改变电子速度、结构周期或数值孔径不能被解释为同一个参数的效应。

\begin{exercise}
在 \(k_z=0\) 的一维光栅里解 \(\omega/v=G_m\) 对周期 \(a\) 的条件，并说明有限长度 \(L\) 将该条件展宽多少。
\end{exercise}
